\begin{table}[!htbp]
\centering
\caption{\textbf{Within-Subject Switching in the Dictator Game}}
\label{tab:within_switching_summary}
\resizebox{0.98\textwidth}{!}{%
\begin{tabular}{lcccccc}
\toprule
 & & \multicolumn{3}{c}{Share switching} & \multicolumn{2}{c}{P(choice switch $\mid$ category)} \\
\cmidrule(lr){3-5} \cmidrule(lr){6-7}
Comparison & $N$ pairs & Choice & Category & Social proximity & Switched & Stayed \\
\midrule
DG-KW vs DG-LT (abstract frame) & 600 & 46.2\% & 22.7\% & 40.5\% & 80.1\% & 36.2\% \\
Market vs Control (DG-KW) & 599 & 73.1\% & 51.6\% & 66.4\% & 82.2\% & 63.4\% \\
Market vs Control (DG-LT) & 600 & 82.3\% & 55.0\% & 62.0\% & 88.5\% & 74.8\% \\
\bottomrule
\end{tabular}
}
\begin{flushleft}
\footnotesize Notes: Each participant makes two dictator-game decisions. In the first row, both decisions use the abstract (control) frame and the game varies (KW vs LT); in the second and third rows, the game is fixed and the frame varies (abstract vs market). ``Choice'' switching is any change in the share sent; ``Category'' switching is a change in the reason-based representation category (Moral, Self-interest, Mutual Benefit/Cooperation, No clear justification); ``Social proximity'' switching is a change in the social-proximity class of the representation. The last two columns report the probability of a choice switch conditional on the representation category switching or staying fixed.
\end{flushleft}
\end{table}
